\subsection{The groups $C_2\times C_{34}$ and $C_{68}$}

\begin{lemma}\label{lem:decic234-finite-sieve}
Put $X=X_1(2,34)$ and let $H\subset J_1(34)(\mathbb F_3)$ be the
reduced rational cuspidal subgroup. If $T$ is an effective divisor
of degree $10$ on $X_{\mathbb F_3}$ whose support is not entirely cuspidal, then
some rational level automorphism $g$ and one of the four rational
maps $f:X\to X_1(34)$ described below satisfy
\[
 [f_*gT-10c_0]\notin H.
\]
Here $c_0$ is a fixed rational cusp on $X_1(34)$.
\end{lemma}
\begin{proof}
The source has genus $49$ and good reduction at $3$.
Its exact noncuspidal closed-point counts in degrees one through ten are
\[
 (0,0,0,36,48,48,144,738,2160,6504),
\]
and its cusp signature is $1^{24},8^3$.
Newton identities for the Frobenius polynomial of
$J_H(68)$ with $H=\langle35\rangle$ independently verify the counts.

For a level structure $(E,P,Q)$, put $R=17Q$, $V=18Q$ and
$(T_1,T_2,T_3)=(P,R,P+R)$.
The $48$ rational level automorphisms are
\[
 (P,Q)\longmapsto(T_i,T_j+aV),\qquad i\ne j,\quad
 i,j\in\{1,2,3\},\quad 1\leq a\leq8.
\]
Their complete orbit counts in degrees four through ten are
\[
 (2,1,2,3,18,45,138).
\]
All $9678$ noncuspidal closed points occur in these $209$ orbits.
The noncuspidal part of $T$ has one support of degree $4,\ldots,10$,
or two supports of degrees $(4,4),(4,5),(4,6),(5,5)$.
We normalize the first support by the level group and retain every
actual closed point for the second. Repetitions are included.
There are $209$ single-support cases and respectively
$72,96,96,48$ mixed cases, giving $521$ tests in total.

The four maps retain $Q$ or $P+Q$, or retain the image of $Q$
after quotienting by $\langle P\rangle$ or
$\langle P+17Q\rangle$. The quotient kernels are disjoint from
$\langle Q\rangle$, so the resulting points have order $34$.
By \cite[Theorem~4.13 and Table~2]{DEHMZ},
\[
 J_1(34)(\Q)=[8760,595680]
\]
is generated by rational cuspidal divisor classes.
The proven finite Jacobian at $3$ has invariants
$[4,4,4,8760,1191360]$, so $H$ has index $128$.
Every finite cusp class belongs to $H$.
Thus cuspidal complements do not change the necessary membership
condition. Every norm multiplicity $d/e$ is included when a
source support of degree $d$ maps to a target support of degree $e$.

Exact class calculations exclude all $521$ cases.
Independent finite-field arithmetic verifies the complete level
orbits, all maps and $842$ normed summands.
Separate calculations reconstruct all $328$ archived target images
using their exact finite and infinite ideal identities.
This proves the assertion.
\end{proof}

\begin{lemma}\label{lem:decic234-short-units}
Let $c_1,\ldots,c_{24}$ be the rational cusps on $X_1(2,34)$.
If a nonconstant rational modular unit has divisor
$\sum_i a_i c_i$, then
\[
 \sum_i a_i^2>60.
\]
\end{lemma}
\begin{proof}
The curve is $\Q$-isomorphic to $X_\Delta(68)$ with
$\Delta=\{\pm1,\pm33\}$.
We compute the full principal-divisor lattice supported on its
$24$ rational cusps, using the integral period map.
Concretely, the lattice consists of the degree-zero integer vectors
whose cusp-period vector belongs to integral homology.
Its rank is $23$. This computes the full lattice, rather than a
subgroup generated by selected modular units.

Exact short-vector enumeration gives no nonzero vector of squared
norm at most $60$. An independent computation reconstructs the
period kernel and verifies the same empty enumeration with PARI.
The integral period matrix and the complete lattice basis are
exported. Hence the stated inequality holds.
\end{proof}

\begin{lemma}\label{lem:decic234-gonality}
The $\Q$-gonality of $X_1(2,34)$ is at least $11$.
\end{lemma}
\begin{proof}
By \cite[Table~1]{NV}, its $\mathbb F_3$-gonality is at least $10$.
Suppose there is a rational map of degree $10$.
Specializing a saturated integral pencil at $3$ gives a
nonconstant map of degree at most $10$ on $X_{\mathbb F_3}$.
The finite gonality bound makes its degree exactly $10$ and
excludes base points.

Every fibre of this finite pencil has a class specializing from a
rational divisor of degree $10$ on $X$. The same is true after
any rational level automorphism. Consequently every pushforward
in Lemma~\ref{lem:decic234-finite-sieve} has normalized class in $H$.
That lemma forces all four rational fibres to be entirely cuspidal.

There are exactly $24$ rational points on $X_{\mathbb F_3}$, and
all are cusps. Each lies in one of these four fibres.
Let $n_1\geq n_2\geq n_3\geq n_4$ be the numbers of distinct
rational cusps in the fibres. Then
\[
 n_1+\cdots+n_4=24,\qquad n_1\geq6,\qquad
 n_2\geq\left\lceil\frac{24-n_1}{3}\right\rceil\geq5.
\]
Since the other cusps have degree eight, neither of the two
largest fibres can contain such a cusp. Both are sums of rational
cusps of total degree $10$.

A positive divisor of degree $10$ supported on $n$ distinct rational
points has squared coefficient norm at most
\[
 B(n)=(11-n)^2+n-1.
\]
The two fibres have disjoint support, and the inequalities above give
\[
 B(n_1)+B(n_2)\leq60.
\]
Their difference is principal over $\mathbb F_3$.
Every rational residue cusp is the reduction of a rational global
cusp: the Tate labels show that Frobenius $3$ generates the global
cusp Galois action. Lift the difference to the corresponding
rational cuspidal divisor in characteristic zero.
Manin--Drinfeld makes its class torsion. Its reduction vanishes,
so torsion injectivity at the odd unramified prime $3$ makes it
principal over $\Q$. This is a nonzero vector in the lattice of
Lemma~\ref{lem:decic234-short-units} with squared norm at most $60$,
a contradiction.
\end{proof}

\begin{proposition}\label{prop:decic234}
No elliptic curve over a decic number field contains
$C_2\times C_{34}$ or $C_{68}$.
\end{proposition}
\begin{proof}
The Jacobian of $X=X_1(2,34)$ has rank zero by
\cite[Theorem~3.1]{DEHMZ}.
Suppose that a proposed level structure gives a noncuspidal trace
divisor $D$ of degree $10$.
Its reduction cannot have noncuspidal support, by
Lemma~\ref{lem:decic234-finite-sieve}, since every indicated
pushforward is a rational divisor class before reduction.
Entirely cuspidal reduction is excluded by
Lemma~\ref{lem:decic-cuspidal-reduction} and
Lemma~\ref{lem:decic234-gonality}.
Thus $C_2\times C_{34}$ is impossible.
Lemma~\ref{lem:decic-isogeny-transfer} also excludes $C_{68}$.
\end{proof}
